% solution source: 2025_SKYMAP_sol.pdf pages 5-11 (OM02: Solution and Marking Scheme)
\textbf{(OM02.1)} It is evident from the constellation lines and their
relative positions that
\[ \text{C1: Virgo (Vir)}, \qquad \text{C2: Corvus (Crv)}. \]
\hfill(1 mark: 0.5 for each constellation)

\textbf{(OM02.2a)} Since the constellation boundaries are parallel to the
grid lines, the grid shown in the figure must be the \textbf{Equatorial
Grid}. Circles of declination are centred around the poles. Hence, the
lines that pass through the red squares are declination ($\delta$) lines.
\hfill(1 mark)

\textbf{(OM02.2b)} Lines emanating from the poles and perpendicular to the
declination circles are RA lines. Hence, the lines that pass through the
blue circles are lines of constant RA ($\alpha$). \hfill(1 mark)

\textbf{(OM02.3)} The constellation of Corvus is towards the South of Virgo.
Hence, the pole of the hemisphere in which Corvus lies is the South pole.
(On the map: the left grid pole is ``S'', the right grid pole is ``N''.)
\hfill(1 mark: 0.5 for each pole)

\textbf{(OM02.4)} The two curves present on the map are the Ecliptic (the
dashed sinusoidal curve, marked with small bars) and the Celestial Equator
(the vertical lines midway between the poles, marked with small circles).
\hfill(2 marks: Ecliptic 1.0; Celestial Equator 1.0, or 0.5 if only one
segment of the Celestial Equator is marked)

\textbf{(OM02.5)} Equinoxes are the intersection points of the equator and
ecliptic. The Autumnal equinox (AE) lies in the constellation of Virgo.
Hence the other intersection point is the Vernal equinox (VE).
\hfill(2 marks: equinoxes at the intersections of CE and Ecliptic 1.0;
both equinoxes correctly identified 1.0)

\textbf{(OM02.6)} From the VE the Sun moves northward, so the Sun's motion
is in the direction of increasing RA. \hfill(1 mark)

\textbf{(OM02.7)} The Vernal Equinox (VE) and Autumnal Equinox (AE) were
identified in part (OM02.5). The Right Ascension of the VE is 0\,h, while
the RA of the AE is 12\,h. From part (OM02.6), the Sun moves northward in
the direction of increasing RA, which allows the RA circles to be marked
accordingly. RA is conventionally measured in hours, minutes, and seconds.
The celestial pole has a declination of $90^\circ$, with northern
declinations positive and southern declinations negative.

The values on the official solution map are: red squares (declination)
$+20^\circ$, $-10^\circ$, $-30^\circ$; blue circles (right ascension)
20\,h, 10\,h, 7\,h.
\hfill(3 marks: 0.5 per correct value with appropriate sign and units;
maximum 1.5 for values incorrect but with appropriate sign and units and
consistent with (OM02.2a) and (OM02.2b))

\textbf{(OM02.8)}
\begin{itemize}
\item In the given list there are 3 zodiacal constellations, Aquarius, Leo
  and Sagittarius, which implies that they should lie on the Ecliptic.
\item The Ecliptic passes through the constellation to the top left which is
  about 2 hours ahead in RA from Virgo, hence it is Sagittarius.
\item From the given list Leo and Orion are the constellations which lie on
  the Celestial Equator. Hence, the constellation on the lower right has to
  be Orion.
\item Out of the remaining constellations, i.e.\ Perseus and Cygnus, the
  constellation of Cygnus lies between RA 20\,h and 22\,h.
\item The remaining box is of Ursa Major, which is not present in the given
  list, so it is marked with a cross ($\times$).
\end{itemize}
\hfill(4 marks: 1.0 for each correct identification, including marking
$\times$ or ``UMa'')

% FIGURE NEEDED: solution overlay of Map-OM02 showing poles S and N, the
% Celestial Equator and Ecliptic marked with symbols, VE and AE, the Sun's
% direction arrow, grid values (+20, -10, -30 deg; 20h, 10h, 7h), and the
% shaded areas labelled Sgr, Cyg, Ori and X (UMa)
% (2025_Map-OM02-Solutions-A3.pdf; also 2025_SKYMAP_sol.pdf pages 6-11)
